% Part: proof-theory
% Chapter: natural-deduction
% Section: introduction

\documentclass[../../../include/open-logic-section]{subfiles}

\begin{document}

\olfileid{pt}{nat}{int}

\olsection{ভূমিকা}

স্বাভাবিক নিষ্পাদন হলো !!{proof} পদ্ধতিগুলির একটি পরিবার।
এখানে প্রতিটি যুক্তিসংযোজক বা পরিমাণসূচকের জন্য দুটি অনুমানবিধি:
একটি সংকেতটি ``আনে'' (!!a{introduction} বিধি), অন্যটি
সেটি ``বাদ দেয়'' (!!a{elimination} বিধি)। যেমন,
\Intro{\lif}-এর সিদ্ধান্তে $!A \lif !B$ রূপের !!{formula}s
থাকে, আর \Elim{\lif}-এর পূর্বধারণায় $!A \lif !B$
রূপের !!{formula}s থাকে।

স্বাভাবিক নিষ্পাদনের প্রথম দুটি সংস্করণ ১৯৩০-এর দশকে
Gerhard Gentzen ও Stanisław Jaśkowski প্রবর্তন করেন।
প্রমাণতাত্ত্বিকেরা প্রধানত Gentzen-এর পদ্ধতি আলোচনা করেন;
Jaśkowski-র পদ্ধতি থেকে শিক্ষায় সর্বাধিক ব্যবহৃত পদ্ধতিগুলি
এসেছে। উভয় পদ্ধতিতেই \emph{অনুমান} দিয়ে শুরু করা যায়—
এমন !!{formula}s যেগুলি !!{prove} করতে হয় না, কিন্তু
অন্য !!{formula}s !!{proving}-এ ব্যবহার করা যায়। গোটা
!!{proof} এমন অনুমানের উপর নির্ভর করতে পারে। অনুমানগুলি
অপ্রমাণিত বলে !!{proof} তার সিদ্ধান্ত (তার
\emph{অন্তিম-!!{formula}}) প্রতিষ্ঠা করে না; কেবল দেখায়,
সিদ্ধান্তটি অনুমানগুলি থেকে অনুসৃত হয়।

কিছু অনুমানবিধির সিদ্ধান্ত আর কিছু অনুমানের উপর নির্ভর করে
না: বিধিগুলি সেই অনুমানগুলি ``অব্যাহতি'' দেয়। যেমন,
\Intro{\lif} বিধি অনুমান~$!A$ থেকে সিদ্ধান্ত~$!B$-র
!!a{proof} নিয়ে $!A \lif !B$ সিদ্ধান্ত দেয়।
$!A \lif !B$-তে শেষ হওয়া !!{proof} আর~$!A$-র
উপর নির্ভর করে না; অনুমান~$!A$ অব্যাহতি পায়।
Gentzen-এর পদ্ধতিতে !!{proof}s হলো !!{formula}s-এর
বৃক্ষ, অনুমানগুলি সর্বোচ্চে থাকা !!{formula}s, আর
\Intro{\lif}-এর মতো বিধির লেবেল জানায় কোন অনুমান
অব্যাহতি পেয়েছে। Jaśkowski-র পদ্ধতিতে কোনো অনুমানের
উপর নির্ভরশীল !!{proof}-এর অংশ আলাদা করে দেখানো হয়,
যেমন উল্লম্ব রেখার পরে সরিয়ে বা বাক্সের ভিতর রেখে। বাক্সে
প্রথম সারিতে থাকে অনুমান~$!A$, শেষ সারিতে শর্তাধীন
সূত্রের অনুবর্তী~$!B$; বাক্সের বাইরে থাকে
\Intro{\lif}-এর সিদ্ধান্ত~$!A \lif !B$, আর ওই অনুমান
অব্যাহতি পায়।
% BN-SRC-542 TARGET-CORRECTION-BEGIN
\par\smallskip\noindent\textnormal{\small\textbf{উৎস-সংশোধন (BN-SRC-542):} বাক্সের প্রথম সারিতে অনুমান $!A$ এবং শেষ সারিতে অনুবর্তী $!B$ থাকে; উৎসের বাক্যটি $!B$-কেই প্রথম ও শেষ সারি বলেছে।}
% BN-SRC-542 TARGET-CORRECTION-END

এই পদ্ধতিগুলি ``স্বাভাবিক'', কারণ বিধিগুলি আমাদের
(অন্তত গণিতবিদদের) স্বাভাবিক যুক্তিপদ্ধতির অনুরূপ।
কাল্পনিক ফল (শর্তাধীন সূত্র) প্রতিষ্ঠা করতে পূর্বাংশ
ধরে নিয়ে তা দিয়ে উত্তরাংশ প্রমাণ করি—এটি
\Intro{\lif}। $!A \lif !B$ এবং তার পূর্বাংশ~$!A$
জানলে উত্তরাংশ~$!B$ সিদ্ধান্ত করি—এটি \Elim{\lif}।
একটি বিচ্ছেদ~$!A \lor !B$ ধরে কোনো সিদ্ধান্ত প্রতিষ্ঠায়
ক্ষেত্রবিচারে প্রমাণ ব্যবহার করি—এটি \Elim{\lor}।
সাধারণ দাবি~$\lforall[x][!A(x)]$ প্রমাণ করতে যথেচ্ছ~$c$-র
জন্য $!A(c)$ দেখাই—এটি \Intro{\lforall}। এভাবে চলতে থাকে।

যেমন, Gentzen-রীতির স্বাভাবিক নিষ্পাদনে
$!A \lif (!A \lor !B)$-র একটি সরল !!{proof}:
\[
\AxiomC{$\Discharge{!A}{x}$}
\RightLabel{\Intro{\lor}}
\UnaryInfC{$!A \lor !B$}
\DischargeRule{\Intro{\lif}}{x}
\UnaryInfC{$!A \lif (!A \lor !B)$}
\DisplayProof
\]
এটি অনুমান~$!A$ দিয়ে শুরু হয়; \Intro{\lor} প্রয়োগে
$!A \lor !B$, তারপর \Intro{\lif} প্রয়োগে
$!A \lif (!A \lor !B)$ পাওয়া যায়। \Intro{\lif}
অনুমান~$!A$-কে অব্যাহতি দেয়; লেবেল~$x$ সেটি জানায়।

প্রমাণতাত্ত্বিকেরা Gentzen-এর মূল !!{formula}s-বৃক্ষ
পদ্ধতি বেশি পছন্দ করেন, যদিও আরেকটি রূপভেদও
প্রমাণতত্ত্বে গুরুত্বপূর্ণ। অনুমান~$\Gamma$ থেকে
$!A$-র !!^a{proof} দেখায় যে $!A$, $\Gamma$ থেকে
অনুসৃত হয়, অর্থাৎ $\Gamma \Entails !A$। স্বাভাবিক
নিষ্পাদনের বিধিগুলিকে এই অনুসরণ-সম্পর্ক সম্পর্কে
বক্তব্য হিসেবেও পড়া যায়। যেমন, \Intro{\lif} বলে,
$\Gamma + !A \Entails !B$ হলে
$\Gamma \Entails !A \lif !B$। তাই বিধিগুলিকে সরাসরি
এমনভাবে গঠন করা স্বাভাবিক, যাতে তাদের পূর্বধারণা ও
সিদ্ধান্তে অনুমানগুলি এবং সেগুলি থেকে অনুসৃত সূত্র—উভয়ের
উপর কাজ করে, অর্থাৎ সিকোয়েন্ট~$\Gamma \Sequent !A$-র
উপর। উপরকার সরল প্রমাণটি এমন পদ্ধতিতে হবে:
\[
\Axiom$x:!A \fCenter !A$
\RightLabel{\Intro{\lor}}
\UnaryInf$x:!A \fCenter !A \lor !B$
\DischargeRule{\Intro{\lif}}{x}
\UnaryInf$\fCenter !A \lif (!A \lor !B)$
\DisplayProof
\]
দেখা যাচ্ছে, বিধিগুলি একই: তারা~$\Sequent$-এর ডান
পাশের !!{formula}-র উপর কাজ করে; বাম পাশের
!!{formula}s শুধু প্রতিটি ধাপে নির্ভরশীল অনুমানগুলি জানায়।
% BN-SRC-543 TARGET-CORRECTION-BEGIN
\par\smallskip\noindent\textnormal{\small\textbf{উৎস-সংশোধন (BN-SRC-543):} $\Intro{\lif}$-এর সিদ্ধান্ত $\Gamma\Entails !B$ নয়, $\Gamma\Entails !A\lif !B$; উৎসের আগের পূর্বধারণা ও প্রদর্শিত সিকোয়েন্টবৃক্ষ অনুযায়ী উত্তরাংশ সংশোধিত।}
% BN-SRC-543 TARGET-CORRECTION-END

!!{introduction}/!!{elimination} বিধিজোড়া একা ধ্রুপদী
যুক্তির জন্য পূর্ণ নয়। $\lfalse$-সম্পর্কিত আরও দুটি
বিধি যোগ করতে হয়। প্রথমটি ``অন্তর্জ্ঞানবাদী অসংগতি
বিধি''~\FalseInt, অর্থাৎ ``বিস্ফোরণ'': $\lfalse$ থেকে
যেকোনো কিছু অনুমান করা যায়। দ্বিতীয়টি পরোক্ষ
প্রমাণের ধ্রুপদী নীতি~\FalseCl{} (অথবা ``ধ্রুপদী
অসংগতি বিধি''): $\lnot !A$ থেকে $\lfalse$ !!{prove}
করা গেলে~$!A$ !!{prove} করা হয়েছে। \FalseCl বাদ দিলে
অন্তর্জ্ঞানবাদী যুক্তির উপযুক্ত পদ্ধতি পাই। দুটি বিধিই
বাদ দিলে ``ন্যূনতম যুক্তি'' পাই।

আমরা ধ্রুপদী ও অন্তর্জ্ঞানবাদী যুক্তির Gentzen-রীতির
স্বাভাবিক নিষ্পাদন আলোচনা করব। এগুলি \Log{N1c} ও
\Log{N1i} পদ্ধতি। \Log{N2c} ও \Log{N2i} হলো সংশ্লিষ্ট
পদ্ধতি, যেখানে আলাদা !!{formula}s~$!A$-র বদলে
সিকোয়েন্ট~$\Gamma \Sequent !A$ ব্যবহৃত হয়।

\end{document}
